% Capítulo 1: archivo incluido desde main.tex mediante % Capítulo 1: archivo incluido desde main.tex mediante % Capítulo 1: archivo incluido desde main.tex mediante % Capítulo 1: archivo incluido desde main.tex mediante \input{chapters/01_introduccion}.
\chapter{Introducción}
\label{chap:introduccion}

La ciberseguridad ha pasado de ser una preocupación estrictamente técnica a
formar parte de la actividad cotidiana de cualquier organización. Las personas
usuarias toman decisiones de seguridad cuando reciben un correo inesperado,
comparten un fichero, eligen una contraseña, notifican una anomalía o utilizan
un servicio en la nube. Sin embargo, la documentación corporativa y las guías
de referencia suelen estar distribuidas entre políticas, procedimientos,
portales y manuales con distintos niveles de detalle. Esta dispersión dificulta
que una persona no especialista encuentre, en el momento oportuno, una
orientación comprensible y coherente con las normas de su organización.

El Marco de Ciberseguridad de NIST propone gestionar el riesgo mediante las
funciones \emph{Govern}, \emph{Identify}, \emph{Protect}, \emph{Detect},
\emph{Respond} y \emph{Recover}; la incorporación explícita de la función de
gobierno en la versión 2.0 subraya que el problema no se limita a desplegar
tecnología, sino que también exige decisiones, responsabilidades y comunicación
organizativa \cite{nistCsf2024}. En paralelo, el informe de panorama de
amenazas de ENISA muestra que el entorno europeo continúa expuesto a amenazas
diversas y cambiantes, por lo que las recomendaciones operativas deben revisarse
con frecuencia y contextualizarse \cite{enisaThreat2025}. Un canal de consulta
de primer nivel no sustituye a un equipo de seguridad, pero puede reducir la
fricción entre una duda habitual y el procedimiento que debe seguirse.

Este Trabajo Fin de Máster aborda ese problema mediante el diseño e
implementación de un agente conversacional especializado en ciberseguridad,
configurado en Microsoft Copilot Studio e integrado en una interfaz web Angular.
El agente se concibe como un sistema de apoyo informativo para consultas
frecuentes en un escenario organizativo simulado: phishing, credenciales,
protección de dispositivos, buenas prácticas, notificación de incidentes y
consulta de políticas. La elección de una interfaz web permite ofrecer un punto
de acceso independiente de la consola de administración del agente, mientras
que Copilot Studio proporciona capacidades de conocimiento, orquestación y
publicación que deben configurarse y gobernarse de forma explícita.

\section{Contexto y planteamiento del problema}
\label{sec:contexto-problema}

Las consultas de ciberseguridad suelen tener una naturaleza situacional. Una
pregunta como ``he recibido un correo con un enlace, ¿qué hago?'' no exige sólo
definir el phishing: requiere proponer pasos prudentes, recordar que no se debe
interactuar con el mensaje y dirigir a la persona al canal de notificación
adecuado. INCIBE define el phishing como una técnica que suplanta a una entidad
legítima para obtener información, realizar cargos o infectar un dispositivo
\cite{incibePhishing}. En un entorno de trabajo, la respuesta útil debe ser
breve, accionable y compatible con la política interna, sin convertir al agente
en un sustituto de la gestión de incidentes.

La literatura y las guías operativas de respuesta a incidentes refuerzan esta
distinción. NIST sitúa la respuesta ante incidentes dentro de la gestión continua
del riesgo y persigue mejorar la preparación, detección, respuesta y recuperación
\cite{nistIncident2025}. Por su parte, la guía de INCIBE para pymes recomienda
definir responsables, canales de notificación, procedimientos, registro de
evidencias y mecanismos de mejora continua \cite{incibeCrisis}. Por tanto, el
agente propuesto no debe ejecutar acciones de contención, investigar sistemas ni
decidir la severidad de un incidente real. Su papel es orientar, pedir sólo el
contexto mínimo imprescindible, advertir de los límites y derivar a los canales
humanos establecidos.

Los modelos de lenguaje son atractivos para este caso de uso porque pueden
interpretar múltiples formulaciones de una misma necesidad y producir una
explicación en lenguaje natural. No obstante, esa flexibilidad no garantiza que
la respuesta sea correcta. El perfil de inteligencia artificial generativa del
AI RMF identifica la \emph{confabulación} como la producción de contenido
erróneo o falso presentado con confianza, y considera también los riesgos de
privacidad, seguridad de la información y configuración humano--IA
\cite{nistGenAi2024}. Una respuesta plausible pero equivocada sobre un posible
incidente puede inducir acciones inadecuadas o una falsa sensación de seguridad.

La estrategia adoptada para mitigar este riesgo es fundamentar las respuestas en
un conjunto acotado de fuentes de conocimiento revisadas. La generación
aumentada por recuperación, conocida como RAG por sus siglas en inglés, combina
la generación del modelo con una memoria documental externa; Lewis et al.
plantean este enfoque para tareas intensivas en conocimiento y señalan su valor
para acceder a información explícita y actualizable \cite{lewis2020rag}. En el
presente trabajo no se pretende entrenar un modelo RAG desde cero. Se emplearán
las capacidades de recuperación y fuentes de conocimiento disponibles en Copilot
Studio, aplicando criterios de selección, versionado y trazabilidad de las
fuentes.

\section{Motivación y oportunidad}
\label{sec:motivacion}

El proyecto se justifica por tres motivos. En primer lugar, la disponibilidad de
una ayuda conversacional puede facilitar la concienciación y orientar a las
personas hacia los recursos existentes sin obligarlas a conocer de antemano la
estructura documental de la organización. En segundo lugar, el caso permite
estudiar de forma práctica el equilibrio entre utilidad, seguridad y control en
un sistema basado en IA generativa. Finalmente, integra dos ámbitos habituales
en un desarrollo de ciberseguridad: la configuración segura de un servicio
gestionado y la construcción de una aplicación web que lo expone a usuarios.

La solución propuesta persigue mejorar la localización y la comprensión de
conocimiento fiable; no persigue automatizar decisiones de alto impacto. Cuando
la información disponible no permita una respuesta suficientemente fundamentada,
el comportamiento deseado será reconocer la incertidumbre y recomendar la
consulta con el equipo responsable. Esta propiedad es especialmente relevante
frente al riesgo de sobreconfianza de los usuarios en las respuestas de modelos
generativos, recogido tanto por NIST como por las guías de seguridad para
aplicaciones con modelos de lenguaje \cite{nistGenAi2024,owaspLlm2026}.

\section{Pregunta de trabajo}
\label{sec:pregunta-trabajo}

La pregunta central que guía el TFM es la siguiente:

\begin{quote}
\emph{¿Cómo puede diseñarse e integrarse un agente conversacional de
ciberseguridad, fundamentado en fuentes reconocidas, que ofrezca orientación
útil y trazable sobre consultas frecuentes sin reemplazar los procesos humanos
de gestión de incidentes?}
\end{quote}

De ella se derivan tres cuestiones operativas:

\begin{enumerate}
  \item ¿Qué casos de uso, fuentes y límites de actuación deben definirse antes
  de configurar el agente?
  \item ¿Qué controles de diseño reducen el riesgo de respuestas no
  fundamentadas, divulgación de información sensible o instrucciones fuera de
  alcance?
  \item ¿Cómo debe evaluarse la calidad de las respuestas de forma repetible,
  separando precisión, cobertura, trazabilidad, consistencia y utilidad?
\end{enumerate}

Estas cuestiones transforman el proyecto en algo más que una demostración de
conversación: el artefacto se diseña junto con un protocolo de evaluación y una
argumentación de seguridad. La metodología elegida se desarrolla en el
Capítulo~\ref{chap:objetivos-metodologia}.

\section{Alcance}
\label{sec:alcance}

El alcance funcional previsto se limita a un escenario organizativo simulado y
a consultas de carácter defensivo y educativo. Las categorías iniciales son las
siguientes:

\begin{itemize}
  \item identificación y comunicación inicial de correos, enlaces o adjuntos
  sospechosos;
  \item gestión de contraseñas, autenticación multifactor y protección de
  cuentas;
  \item pautas básicas para dispositivos, actualizaciones, copias de seguridad y
  uso seguro de servicios corporativos;
  \item orientación inicial ante señales de incidente y derivación al canal de
  notificación; y
  \item explicación de políticas corporativas simuladas y de fuentes públicas
  seleccionadas.
\end{itemize}

La base de conocimiento se compondrá de documentos o sitios oficiales de NIST,
ENISA, INCIBE y Microsoft Learn, complementados, si procede, por una política
corporativa ficticia creada expresamente para las pruebas. Cada fuente deberá
tener propietario, URL o identificador, fecha de revisión, alcance temático y
estado de aprobación. Las fuentes públicas no se tratarán como evidencia
inmutable: deberán revisarse antes de cada versión evaluada, ya que las guías,
las plataformas y el panorama de amenazas evolucionan.

Quedan fuera del alcance los siguientes aspectos:

\begin{itemize}
  \item la monitorización de sistemas, el análisis forense, la ejecución de
  comandos, la respuesta automática ante incidentes o el acceso a herramientas
  de seguridad;
  \item la recepción o tratamiento intencionado de credenciales, secretos,
  registros de producción, datos personales reales o información clasificada;
  \item las instrucciones ofensivas que puedan facilitar intrusión, evasión,
  malware o exfiltración;
  \item la sustitución de un SOC, de un CSIRT, de los responsables de protección
  de datos o de asesoramiento jurídico; y
  \item la afirmación de conformidad legal o certificación de seguridad de una
  organización real.
\end{itemize}

Esta delimitación es deliberada. El agente debe indicar que, ante un riesgo
inmediato o una posible brecha, la persona usuaria ha de seguir el procedimiento
organizativo y contactar con el equipo responsable. Si no existe tal
procedimiento en el entorno de demostración, el texto de derivación se marcará
como un elemento configurable, no como una instrucción universal.

\section{Consideraciones de seguridad, privacidad y gobierno}
\label{sec:consideraciones-intro}

El sistema integra dos superficies que requieren atención: el comportamiento del
agente y la aplicación web que lo presenta. Entre los riesgos característicos de
las aplicaciones con modelos de lenguaje se encuentran la inyección de
instrucciones, la divulgación de información sensible, la contaminación de datos
y la dependencia excesiva de las salidas del modelo \cite{owaspLlm2026}. En
consecuencia, las instrucciones del agente deben priorizar fuentes aprobadas,
prohibir la solicitud de secretos, limitar la respuesta a su dominio y disponer
de rutas de rechazo o escalado.

En el plano web, Angular ofrece mecanismos de protección frente a XSS al tratar
los valores de las plantillas como no confiables y recomienda mantener el
framework actualizado, evitar APIs de riesgo y complementar las defensas con una
política de seguridad de contenido y \emph{Trusted Types} cuando sea aplicable
\cite{angularSecurity}. Estas medidas no sustituyen a la autenticación,
autorización, protección de secretos ni validación del lado servidor; forman
parte de una defensa en profundidad que se concretará durante la implementación.

El Reglamento General de Protección de Datos exige que el tratamiento de datos
personales se realice con medidas técnicas y organizativas adecuadas al riesgo
\cite{gdpr}. Además, el Reglamento europeo de IA establece un marco armonizado
para el desarrollo y uso de sistemas de IA orientado a una IA fiable y centrada
en las personas \cite{aiAct}. Este TFM no emite un dictamen de cumplimiento.
Sí incorpora esos marcos como criterios de diseño: minimización de datos,
transparencia sobre las limitaciones, supervisión humana y documentación de las
decisiones de configuración.

\section{Estructura de la memoria}
\label{sec:estructura-memoria}

El Capítulo~\ref{chap:objetivos-metodologia} define los objetivos, requisitos,
metodología y protocolo de evaluación. El Capítulo~\ref{chap:estado-arte}
revisa los fundamentos de agentes con recuperación de conocimiento, Microsoft
Copilot Studio, seguridad de aplicaciones web y gobierno de IA. Los capítulos
posteriores describirán la arquitectura implementada, la configuración del
agente, la integración con Angular y las decisiones de seguridad. A continuación
se documentará la validación mediante casos de prueba sintéticos y sus resultados
reproducibles. Por último, las conclusiones contrastarán los objetivos con la
evidencia obtenida y propondrán mejoras para un posible despliegue organizativo.

.
\chapter{Introducción}
\label{chap:introduccion}

La ciberseguridad ha pasado de ser una preocupación estrictamente técnica a
formar parte de la actividad cotidiana de cualquier organización. Las personas
usuarias toman decisiones de seguridad cuando reciben un correo inesperado,
comparten un fichero, eligen una contraseña, notifican una anomalía o utilizan
un servicio en la nube. Sin embargo, la documentación corporativa y las guías
de referencia suelen estar distribuidas entre políticas, procedimientos,
portales y manuales con distintos niveles de detalle. Esta dispersión dificulta
que una persona no especialista encuentre, en el momento oportuno, una
orientación comprensible y coherente con las normas de su organización.

El Marco de Ciberseguridad de NIST propone gestionar el riesgo mediante las
funciones \emph{Govern}, \emph{Identify}, \emph{Protect}, \emph{Detect},
\emph{Respond} y \emph{Recover}; la incorporación explícita de la función de
gobierno en la versión 2.0 subraya que el problema no se limita a desplegar
tecnología, sino que también exige decisiones, responsabilidades y comunicación
organizativa \cite{nistCsf2024}. En paralelo, el informe de panorama de
amenazas de ENISA muestra que el entorno europeo continúa expuesto a amenazas
diversas y cambiantes, por lo que las recomendaciones operativas deben revisarse
con frecuencia y contextualizarse \cite{enisaThreat2025}. Un canal de consulta
de primer nivel no sustituye a un equipo de seguridad, pero puede reducir la
fricción entre una duda habitual y el procedimiento que debe seguirse.

Este Trabajo Fin de Máster aborda ese problema mediante el diseño e
implementación de un agente conversacional especializado en ciberseguridad,
configurado en Microsoft Copilot Studio e integrado en una interfaz web Angular.
El agente se concibe como un sistema de apoyo informativo para consultas
frecuentes en un escenario organizativo simulado: phishing, credenciales,
protección de dispositivos, buenas prácticas, notificación de incidentes y
consulta de políticas. La elección de una interfaz web permite ofrecer un punto
de acceso independiente de la consola de administración del agente, mientras
que Copilot Studio proporciona capacidades de conocimiento, orquestación y
publicación que deben configurarse y gobernarse de forma explícita.

\section{Contexto y planteamiento del problema}
\label{sec:contexto-problema}

Las consultas de ciberseguridad suelen tener una naturaleza situacional. Una
pregunta como ``he recibido un correo con un enlace, ¿qué hago?'' no exige sólo
definir el phishing: requiere proponer pasos prudentes, recordar que no se debe
interactuar con el mensaje y dirigir a la persona al canal de notificación
adecuado. INCIBE define el phishing como una técnica que suplanta a una entidad
legítima para obtener información, realizar cargos o infectar un dispositivo
\cite{incibePhishing}. En un entorno de trabajo, la respuesta útil debe ser
breve, accionable y compatible con la política interna, sin convertir al agente
en un sustituto de la gestión de incidentes.

La literatura y las guías operativas de respuesta a incidentes refuerzan esta
distinción. NIST sitúa la respuesta ante incidentes dentro de la gestión continua
del riesgo y persigue mejorar la preparación, detección, respuesta y recuperación
\cite{nistIncident2025}. Por su parte, la guía de INCIBE para pymes recomienda
definir responsables, canales de notificación, procedimientos, registro de
evidencias y mecanismos de mejora continua \cite{incibeCrisis}. Por tanto, el
agente propuesto no debe ejecutar acciones de contención, investigar sistemas ni
decidir la severidad de un incidente real. Su papel es orientar, pedir sólo el
contexto mínimo imprescindible, advertir de los límites y derivar a los canales
humanos establecidos.

Los modelos de lenguaje son atractivos para este caso de uso porque pueden
interpretar múltiples formulaciones de una misma necesidad y producir una
explicación en lenguaje natural. No obstante, esa flexibilidad no garantiza que
la respuesta sea correcta. El perfil de inteligencia artificial generativa del
AI RMF identifica la \emph{confabulación} como la producción de contenido
erróneo o falso presentado con confianza, y considera también los riesgos de
privacidad, seguridad de la información y configuración humano--IA
\cite{nistGenAi2024}. Una respuesta plausible pero equivocada sobre un posible
incidente puede inducir acciones inadecuadas o una falsa sensación de seguridad.

La estrategia adoptada para mitigar este riesgo es fundamentar las respuestas en
un conjunto acotado de fuentes de conocimiento revisadas. La generación
aumentada por recuperación, conocida como RAG por sus siglas en inglés, combina
la generación del modelo con una memoria documental externa; Lewis et al.
plantean este enfoque para tareas intensivas en conocimiento y señalan su valor
para acceder a información explícita y actualizable \cite{lewis2020rag}. En el
presente trabajo no se pretende entrenar un modelo RAG desde cero. Se emplearán
las capacidades de recuperación y fuentes de conocimiento disponibles en Copilot
Studio, aplicando criterios de selección, versionado y trazabilidad de las
fuentes.

\section{Motivación y oportunidad}
\label{sec:motivacion}

El proyecto se justifica por tres motivos. En primer lugar, la disponibilidad de
una ayuda conversacional puede facilitar la concienciación y orientar a las
personas hacia los recursos existentes sin obligarlas a conocer de antemano la
estructura documental de la organización. En segundo lugar, el caso permite
estudiar de forma práctica el equilibrio entre utilidad, seguridad y control en
un sistema basado en IA generativa. Finalmente, integra dos ámbitos habituales
en un desarrollo de ciberseguridad: la configuración segura de un servicio
gestionado y la construcción de una aplicación web que lo expone a usuarios.

La solución propuesta persigue mejorar la localización y la comprensión de
conocimiento fiable; no persigue automatizar decisiones de alto impacto. Cuando
la información disponible no permita una respuesta suficientemente fundamentada,
el comportamiento deseado será reconocer la incertidumbre y recomendar la
consulta con el equipo responsable. Esta propiedad es especialmente relevante
frente al riesgo de sobreconfianza de los usuarios en las respuestas de modelos
generativos, recogido tanto por NIST como por las guías de seguridad para
aplicaciones con modelos de lenguaje \cite{nistGenAi2024,owaspLlm2026}.

\section{Pregunta de trabajo}
\label{sec:pregunta-trabajo}

La pregunta central que guía el TFM es la siguiente:

\begin{quote}
\emph{¿Cómo puede diseñarse e integrarse un agente conversacional de
ciberseguridad, fundamentado en fuentes reconocidas, que ofrezca orientación
útil y trazable sobre consultas frecuentes sin reemplazar los procesos humanos
de gestión de incidentes?}
\end{quote}

De ella se derivan tres cuestiones operativas:

\begin{enumerate}
  \item ¿Qué casos de uso, fuentes y límites de actuación deben definirse antes
  de configurar el agente?
  \item ¿Qué controles de diseño reducen el riesgo de respuestas no
  fundamentadas, divulgación de información sensible o instrucciones fuera de
  alcance?
  \item ¿Cómo debe evaluarse la calidad de las respuestas de forma repetible,
  separando precisión, cobertura, trazabilidad, consistencia y utilidad?
\end{enumerate}

Estas cuestiones transforman el proyecto en algo más que una demostración de
conversación: el artefacto se diseña junto con un protocolo de evaluación y una
argumentación de seguridad. La metodología elegida se desarrolla en el
Capítulo~\ref{chap:objetivos-metodologia}.

\section{Alcance}
\label{sec:alcance}

El alcance funcional previsto se limita a un escenario organizativo simulado y
a consultas de carácter defensivo y educativo. Las categorías iniciales son las
siguientes:

\begin{itemize}
  \item identificación y comunicación inicial de correos, enlaces o adjuntos
  sospechosos;
  \item gestión de contraseñas, autenticación multifactor y protección de
  cuentas;
  \item pautas básicas para dispositivos, actualizaciones, copias de seguridad y
  uso seguro de servicios corporativos;
  \item orientación inicial ante señales de incidente y derivación al canal de
  notificación; y
  \item explicación de políticas corporativas simuladas y de fuentes públicas
  seleccionadas.
\end{itemize}

La base de conocimiento se compondrá de documentos o sitios oficiales de NIST,
ENISA, INCIBE y Microsoft Learn, complementados, si procede, por una política
corporativa ficticia creada expresamente para las pruebas. Cada fuente deberá
tener propietario, URL o identificador, fecha de revisión, alcance temático y
estado de aprobación. Las fuentes públicas no se tratarán como evidencia
inmutable: deberán revisarse antes de cada versión evaluada, ya que las guías,
las plataformas y el panorama de amenazas evolucionan.

Quedan fuera del alcance los siguientes aspectos:

\begin{itemize}
  \item la monitorización de sistemas, el análisis forense, la ejecución de
  comandos, la respuesta automática ante incidentes o el acceso a herramientas
  de seguridad;
  \item la recepción o tratamiento intencionado de credenciales, secretos,
  registros de producción, datos personales reales o información clasificada;
  \item las instrucciones ofensivas que puedan facilitar intrusión, evasión,
  malware o exfiltración;
  \item la sustitución de un SOC, de un CSIRT, de los responsables de protección
  de datos o de asesoramiento jurídico; y
  \item la afirmación de conformidad legal o certificación de seguridad de una
  organización real.
\end{itemize}

Esta delimitación es deliberada. El agente debe indicar que, ante un riesgo
inmediato o una posible brecha, la persona usuaria ha de seguir el procedimiento
organizativo y contactar con el equipo responsable. Si no existe tal
procedimiento en el entorno de demostración, el texto de derivación se marcará
como un elemento configurable, no como una instrucción universal.

\section{Consideraciones de seguridad, privacidad y gobierno}
\label{sec:consideraciones-intro}

El sistema integra dos superficies que requieren atención: el comportamiento del
agente y la aplicación web que lo presenta. Entre los riesgos característicos de
las aplicaciones con modelos de lenguaje se encuentran la inyección de
instrucciones, la divulgación de información sensible, la contaminación de datos
y la dependencia excesiva de las salidas del modelo \cite{owaspLlm2026}. En
consecuencia, las instrucciones del agente deben priorizar fuentes aprobadas,
prohibir la solicitud de secretos, limitar la respuesta a su dominio y disponer
de rutas de rechazo o escalado.

En el plano web, Angular ofrece mecanismos de protección frente a XSS al tratar
los valores de las plantillas como no confiables y recomienda mantener el
framework actualizado, evitar APIs de riesgo y complementar las defensas con una
política de seguridad de contenido y \emph{Trusted Types} cuando sea aplicable
\cite{angularSecurity}. Estas medidas no sustituyen a la autenticación,
autorización, protección de secretos ni validación del lado servidor; forman
parte de una defensa en profundidad que se concretará durante la implementación.

El Reglamento General de Protección de Datos exige que el tratamiento de datos
personales se realice con medidas técnicas y organizativas adecuadas al riesgo
\cite{gdpr}. Además, el Reglamento europeo de IA establece un marco armonizado
para el desarrollo y uso de sistemas de IA orientado a una IA fiable y centrada
en las personas \cite{aiAct}. Este TFM no emite un dictamen de cumplimiento.
Sí incorpora esos marcos como criterios de diseño: minimización de datos,
transparencia sobre las limitaciones, supervisión humana y documentación de las
decisiones de configuración.

\section{Estructura de la memoria}
\label{sec:estructura-memoria}

El Capítulo~\ref{chap:objetivos-metodologia} define los objetivos, requisitos,
metodología y protocolo de evaluación. El Capítulo~\ref{chap:estado-arte}
revisa los fundamentos de agentes con recuperación de conocimiento, Microsoft
Copilot Studio, seguridad de aplicaciones web y gobierno de IA. Los capítulos
posteriores describirán la arquitectura implementada, la configuración del
agente, la integración con Angular y las decisiones de seguridad. A continuación
se documentará la validación mediante casos de prueba sintéticos y sus resultados
reproducibles. Por último, las conclusiones contrastarán los objetivos con la
evidencia obtenida y propondrán mejoras para un posible despliegue organizativo.

.
\chapter{Introducción}
\label{chap:introduccion}

La ciberseguridad ha pasado de ser una preocupación estrictamente técnica a
formar parte de la actividad cotidiana de cualquier organización. Las personas
usuarias toman decisiones de seguridad cuando reciben un correo inesperado,
comparten un fichero, eligen una contraseña, notifican una anomalía o utilizan
un servicio en la nube. Sin embargo, la documentación corporativa y las guías
de referencia suelen estar distribuidas entre políticas, procedimientos,
portales y manuales con distintos niveles de detalle. Esta dispersión dificulta
que una persona no especialista encuentre, en el momento oportuno, una
orientación comprensible y coherente con las normas de su organización.

El Marco de Ciberseguridad de NIST propone gestionar el riesgo mediante las
funciones \emph{Govern}, \emph{Identify}, \emph{Protect}, \emph{Detect},
\emph{Respond} y \emph{Recover}; la incorporación explícita de la función de
gobierno en la versión 2.0 subraya que el problema no se limita a desplegar
tecnología, sino que también exige decisiones, responsabilidades y comunicación
organizativa \cite{nistCsf2024}. En paralelo, el informe de panorama de
amenazas de ENISA muestra que el entorno europeo continúa expuesto a amenazas
diversas y cambiantes, por lo que las recomendaciones operativas deben revisarse
con frecuencia y contextualizarse \cite{enisaThreat2025}. Un canal de consulta
de primer nivel no sustituye a un equipo de seguridad, pero puede reducir la
fricción entre una duda habitual y el procedimiento que debe seguirse.

Este Trabajo Fin de Máster aborda ese problema mediante el diseño e
implementación de un agente conversacional especializado en ciberseguridad,
configurado en Microsoft Copilot Studio e integrado en una interfaz web Angular.
El agente se concibe como un sistema de apoyo informativo para consultas
frecuentes en un escenario organizativo simulado: phishing, credenciales,
protección de dispositivos, buenas prácticas, notificación de incidentes y
consulta de políticas. La elección de una interfaz web permite ofrecer un punto
de acceso independiente de la consola de administración del agente, mientras
que Copilot Studio proporciona capacidades de conocimiento, orquestación y
publicación que deben configurarse y gobernarse de forma explícita.

\section{Contexto y planteamiento del problema}
\label{sec:contexto-problema}

Las consultas de ciberseguridad suelen tener una naturaleza situacional. Una
pregunta como ``he recibido un correo con un enlace, ¿qué hago?'' no exige sólo
definir el phishing: requiere proponer pasos prudentes, recordar que no se debe
interactuar con el mensaje y dirigir a la persona al canal de notificación
adecuado. INCIBE define el phishing como una técnica que suplanta a una entidad
legítima para obtener información, realizar cargos o infectar un dispositivo
\cite{incibePhishing}. En un entorno de trabajo, la respuesta útil debe ser
breve, accionable y compatible con la política interna, sin convertir al agente
en un sustituto de la gestión de incidentes.

La literatura y las guías operativas de respuesta a incidentes refuerzan esta
distinción. NIST sitúa la respuesta ante incidentes dentro de la gestión continua
del riesgo y persigue mejorar la preparación, detección, respuesta y recuperación
\cite{nistIncident2025}. Por su parte, la guía de INCIBE para pymes recomienda
definir responsables, canales de notificación, procedimientos, registro de
evidencias y mecanismos de mejora continua \cite{incibeCrisis}. Por tanto, el
agente propuesto no debe ejecutar acciones de contención, investigar sistemas ni
decidir la severidad de un incidente real. Su papel es orientar, pedir sólo el
contexto mínimo imprescindible, advertir de los límites y derivar a los canales
humanos establecidos.

Los modelos de lenguaje son atractivos para este caso de uso porque pueden
interpretar múltiples formulaciones de una misma necesidad y producir una
explicación en lenguaje natural. No obstante, esa flexibilidad no garantiza que
la respuesta sea correcta. El perfil de inteligencia artificial generativa del
AI RMF identifica la \emph{confabulación} como la producción de contenido
erróneo o falso presentado con confianza, y considera también los riesgos de
privacidad, seguridad de la información y configuración humano--IA
\cite{nistGenAi2024}. Una respuesta plausible pero equivocada sobre un posible
incidente puede inducir acciones inadecuadas o una falsa sensación de seguridad.

La estrategia adoptada para mitigar este riesgo es fundamentar las respuestas en
un conjunto acotado de fuentes de conocimiento revisadas. La generación
aumentada por recuperación, conocida como RAG por sus siglas en inglés, combina
la generación del modelo con una memoria documental externa; Lewis et al.
plantean este enfoque para tareas intensivas en conocimiento y señalan su valor
para acceder a información explícita y actualizable \cite{lewis2020rag}. En el
presente trabajo no se pretende entrenar un modelo RAG desde cero. Se emplearán
las capacidades de recuperación y fuentes de conocimiento disponibles en Copilot
Studio, aplicando criterios de selección, versionado y trazabilidad de las
fuentes.

\section{Motivación y oportunidad}
\label{sec:motivacion}

El proyecto se justifica por tres motivos. En primer lugar, la disponibilidad de
una ayuda conversacional puede facilitar la concienciación y orientar a las
personas hacia los recursos existentes sin obligarlas a conocer de antemano la
estructura documental de la organización. En segundo lugar, el caso permite
estudiar de forma práctica el equilibrio entre utilidad, seguridad y control en
un sistema basado en IA generativa. Finalmente, integra dos ámbitos habituales
en un desarrollo de ciberseguridad: la configuración segura de un servicio
gestionado y la construcción de una aplicación web que lo expone a usuarios.

La solución propuesta persigue mejorar la localización y la comprensión de
conocimiento fiable; no persigue automatizar decisiones de alto impacto. Cuando
la información disponible no permita una respuesta suficientemente fundamentada,
el comportamiento deseado será reconocer la incertidumbre y recomendar la
consulta con el equipo responsable. Esta propiedad es especialmente relevante
frente al riesgo de sobreconfianza de los usuarios en las respuestas de modelos
generativos, recogido tanto por NIST como por las guías de seguridad para
aplicaciones con modelos de lenguaje \cite{nistGenAi2024,owaspLlm2026}.

\section{Pregunta de trabajo}
\label{sec:pregunta-trabajo}

La pregunta central que guía el TFM es la siguiente:

\begin{quote}
\emph{¿Cómo puede diseñarse e integrarse un agente conversacional de
ciberseguridad, fundamentado en fuentes reconocidas, que ofrezca orientación
útil y trazable sobre consultas frecuentes sin reemplazar los procesos humanos
de gestión de incidentes?}
\end{quote}

De ella se derivan tres cuestiones operativas:

\begin{enumerate}
  \item ¿Qué casos de uso, fuentes y límites de actuación deben definirse antes
  de configurar el agente?
  \item ¿Qué controles de diseño reducen el riesgo de respuestas no
  fundamentadas, divulgación de información sensible o instrucciones fuera de
  alcance?
  \item ¿Cómo debe evaluarse la calidad de las respuestas de forma repetible,
  separando precisión, cobertura, trazabilidad, consistencia y utilidad?
\end{enumerate}

Estas cuestiones transforman el proyecto en algo más que una demostración de
conversación: el artefacto se diseña junto con un protocolo de evaluación y una
argumentación de seguridad. La metodología elegida se desarrolla en el
Capítulo~\ref{chap:objetivos-metodologia}.

\section{Alcance}
\label{sec:alcance}

El alcance funcional previsto se limita a un escenario organizativo simulado y
a consultas de carácter defensivo y educativo. Las categorías iniciales son las
siguientes:

\begin{itemize}
  \item identificación y comunicación inicial de correos, enlaces o adjuntos
  sospechosos;
  \item gestión de contraseñas, autenticación multifactor y protección de
  cuentas;
  \item pautas básicas para dispositivos, actualizaciones, copias de seguridad y
  uso seguro de servicios corporativos;
  \item orientación inicial ante señales de incidente y derivación al canal de
  notificación; y
  \item explicación de políticas corporativas simuladas y de fuentes públicas
  seleccionadas.
\end{itemize}

La base de conocimiento se compondrá de documentos o sitios oficiales de NIST,
ENISA, INCIBE y Microsoft Learn, complementados, si procede, por una política
corporativa ficticia creada expresamente para las pruebas. Cada fuente deberá
tener propietario, URL o identificador, fecha de revisión, alcance temático y
estado de aprobación. Las fuentes públicas no se tratarán como evidencia
inmutable: deberán revisarse antes de cada versión evaluada, ya que las guías,
las plataformas y el panorama de amenazas evolucionan.

Quedan fuera del alcance los siguientes aspectos:

\begin{itemize}
  \item la monitorización de sistemas, el análisis forense, la ejecución de
  comandos, la respuesta automática ante incidentes o el acceso a herramientas
  de seguridad;
  \item la recepción o tratamiento intencionado de credenciales, secretos,
  registros de producción, datos personales reales o información clasificada;
  \item las instrucciones ofensivas que puedan facilitar intrusión, evasión,
  malware o exfiltración;
  \item la sustitución de un SOC, de un CSIRT, de los responsables de protección
  de datos o de asesoramiento jurídico; y
  \item la afirmación de conformidad legal o certificación de seguridad de una
  organización real.
\end{itemize}

Esta delimitación es deliberada. El agente debe indicar que, ante un riesgo
inmediato o una posible brecha, la persona usuaria ha de seguir el procedimiento
organizativo y contactar con el equipo responsable. Si no existe tal
procedimiento en el entorno de demostración, el texto de derivación se marcará
como un elemento configurable, no como una instrucción universal.

\section{Consideraciones de seguridad, privacidad y gobierno}
\label{sec:consideraciones-intro}

El sistema integra dos superficies que requieren atención: el comportamiento del
agente y la aplicación web que lo presenta. Entre los riesgos característicos de
las aplicaciones con modelos de lenguaje se encuentran la inyección de
instrucciones, la divulgación de información sensible, la contaminación de datos
y la dependencia excesiva de las salidas del modelo \cite{owaspLlm2026}. En
consecuencia, las instrucciones del agente deben priorizar fuentes aprobadas,
prohibir la solicitud de secretos, limitar la respuesta a su dominio y disponer
de rutas de rechazo o escalado.

En el plano web, Angular ofrece mecanismos de protección frente a XSS al tratar
los valores de las plantillas como no confiables y recomienda mantener el
framework actualizado, evitar APIs de riesgo y complementar las defensas con una
política de seguridad de contenido y \emph{Trusted Types} cuando sea aplicable
\cite{angularSecurity}. Estas medidas no sustituyen a la autenticación,
autorización, protección de secretos ni validación del lado servidor; forman
parte de una defensa en profundidad que se concretará durante la implementación.

El Reglamento General de Protección de Datos exige que el tratamiento de datos
personales se realice con medidas técnicas y organizativas adecuadas al riesgo
\cite{gdpr}. Además, el Reglamento europeo de IA establece un marco armonizado
para el desarrollo y uso de sistemas de IA orientado a una IA fiable y centrada
en las personas \cite{aiAct}. Este TFM no emite un dictamen de cumplimiento.
Sí incorpora esos marcos como criterios de diseño: minimización de datos,
transparencia sobre las limitaciones, supervisión humana y documentación de las
decisiones de configuración.

\section{Estructura de la memoria}
\label{sec:estructura-memoria}

El Capítulo~\ref{chap:objetivos-metodologia} define los objetivos, requisitos,
metodología y protocolo de evaluación. El Capítulo~\ref{chap:estado-arte}
revisa los fundamentos de agentes con recuperación de conocimiento, Microsoft
Copilot Studio, seguridad de aplicaciones web y gobierno de IA. Los capítulos
posteriores describirán la arquitectura implementada, la configuración del
agente, la integración con Angular y las decisiones de seguridad. A continuación
se documentará la validación mediante casos de prueba sintéticos y sus resultados
reproducibles. Por último, las conclusiones contrastarán los objetivos con la
evidencia obtenida y propondrán mejoras para un posible despliegue organizativo.

.
\chapter{Introducción}
\label{chap:introduccion}

La ciberseguridad ha pasado de ser una preocupación estrictamente técnica a
formar parte de la actividad cotidiana de cualquier organización. Las personas
usuarias toman decisiones de seguridad cuando reciben un correo inesperado,
comparten un fichero, eligen una contraseña, notifican una anomalía o utilizan
un servicio en la nube. Sin embargo, la documentación corporativa y las guías
de referencia suelen estar distribuidas entre políticas, procedimientos,
portales y manuales con distintos niveles de detalle. Esta dispersión dificulta
que una persona no especialista encuentre, en el momento oportuno, una
orientación comprensible y coherente con las normas de su organización.

El Marco de Ciberseguridad de NIST propone gestionar el riesgo mediante las
funciones \emph{Govern}, \emph{Identify}, \emph{Protect}, \emph{Detect},
\emph{Respond} y \emph{Recover}; la incorporación explícita de la función de
gobierno en la versión 2.0 subraya que el problema no se limita a desplegar
tecnología, sino que también exige decisiones, responsabilidades y comunicación
organizativa \cite{nistCsf2024}. En paralelo, el informe de panorama de
amenazas de ENISA muestra que el entorno europeo continúa expuesto a amenazas
diversas y cambiantes, por lo que las recomendaciones operativas deben revisarse
con frecuencia y contextualizarse \cite{enisaThreat2025}. Un canal de consulta
de primer nivel no sustituye a un equipo de seguridad, pero puede reducir la
fricción entre una duda habitual y el procedimiento que debe seguirse.

Este Trabajo Fin de Máster aborda ese problema mediante el diseño e
implementación de un agente conversacional especializado en ciberseguridad,
configurado en Microsoft Copilot Studio e integrado en una interfaz web Angular.
El agente se concibe como un sistema de apoyo informativo para consultas
frecuentes en un escenario organizativo simulado: phishing, credenciales,
protección de dispositivos, buenas prácticas, notificación de incidentes y
consulta de políticas. La elección de una interfaz web permite ofrecer un punto
de acceso independiente de la consola de administración del agente, mientras
que Copilot Studio proporciona capacidades de conocimiento, orquestación y
publicación que deben configurarse y gobernarse de forma explícita.

\section{Contexto y planteamiento del problema}
\label{sec:contexto-problema}

Las consultas de ciberseguridad suelen tener una naturaleza situacional. Una
pregunta como ``he recibido un correo con un enlace, ¿qué hago?'' no exige sólo
definir el phishing: requiere proponer pasos prudentes, recordar que no se debe
interactuar con el mensaje y dirigir a la persona al canal de notificación
adecuado. INCIBE define el phishing como una técnica que suplanta a una entidad
legítima para obtener información, realizar cargos o infectar un dispositivo
\cite{incibePhishing}. En un entorno de trabajo, la respuesta útil debe ser
breve, accionable y compatible con la política interna, sin convertir al agente
en un sustituto de la gestión de incidentes.

La literatura y las guías operativas de respuesta a incidentes refuerzan esta
distinción. NIST sitúa la respuesta ante incidentes dentro de la gestión continua
del riesgo y persigue mejorar la preparación, detección, respuesta y recuperación
\cite{nistIncident2025}. Por su parte, la guía de INCIBE para pymes recomienda
definir responsables, canales de notificación, procedimientos, registro de
evidencias y mecanismos de mejora continua \cite{incibeCrisis}. Por tanto, el
agente propuesto no debe ejecutar acciones de contención, investigar sistemas ni
decidir la severidad de un incidente real. Su papel es orientar, pedir sólo el
contexto mínimo imprescindible, advertir de los límites y derivar a los canales
humanos establecidos.

Los modelos de lenguaje son atractivos para este caso de uso porque pueden
interpretar múltiples formulaciones de una misma necesidad y producir una
explicación en lenguaje natural. No obstante, esa flexibilidad no garantiza que
la respuesta sea correcta. El perfil de inteligencia artificial generativa del
AI RMF identifica la \emph{confabulación} como la producción de contenido
erróneo o falso presentado con confianza, y considera también los riesgos de
privacidad, seguridad de la información y configuración humano--IA
\cite{nistGenAi2024}. Una respuesta plausible pero equivocada sobre un posible
incidente puede inducir acciones inadecuadas o una falsa sensación de seguridad.

La estrategia adoptada para mitigar este riesgo es fundamentar las respuestas en
un conjunto acotado de fuentes de conocimiento revisadas. La generación
aumentada por recuperación, conocida como RAG por sus siglas en inglés, combina
la generación del modelo con una memoria documental externa; Lewis et al.
plantean este enfoque para tareas intensivas en conocimiento y señalan su valor
para acceder a información explícita y actualizable \cite{lewis2020rag}. En el
presente trabajo no se pretende entrenar un modelo RAG desde cero. Se emplearán
las capacidades de recuperación y fuentes de conocimiento disponibles en Copilot
Studio, aplicando criterios de selección, versionado y trazabilidad de las
fuentes.

\section{Motivación y oportunidad}
\label{sec:motivacion}

El proyecto se justifica por tres motivos. En primer lugar, la disponibilidad de
una ayuda conversacional puede facilitar la concienciación y orientar a las
personas hacia los recursos existentes sin obligarlas a conocer de antemano la
estructura documental de la organización. En segundo lugar, el caso permite
estudiar de forma práctica el equilibrio entre utilidad, seguridad y control en
un sistema basado en IA generativa. Finalmente, integra dos ámbitos habituales
en un desarrollo de ciberseguridad: la configuración segura de un servicio
gestionado y la construcción de una aplicación web que lo expone a usuarios.

La solución propuesta persigue mejorar la localización y la comprensión de
conocimiento fiable; no persigue automatizar decisiones de alto impacto. Cuando
la información disponible no permita una respuesta suficientemente fundamentada,
el comportamiento deseado será reconocer la incertidumbre y recomendar la
consulta con el equipo responsable. Esta propiedad es especialmente relevante
frente al riesgo de sobreconfianza de los usuarios en las respuestas de modelos
generativos, recogido tanto por NIST como por las guías de seguridad para
aplicaciones con modelos de lenguaje \cite{nistGenAi2024,owaspLlm2026}.

\section{Pregunta de trabajo}
\label{sec:pregunta-trabajo}

La pregunta central que guía el TFM es la siguiente:

\begin{quote}
\emph{¿Cómo puede diseñarse e integrarse un agente conversacional de
ciberseguridad, fundamentado en fuentes reconocidas, que ofrezca orientación
útil y trazable sobre consultas frecuentes sin reemplazar los procesos humanos
de gestión de incidentes?}
\end{quote}

De ella se derivan tres cuestiones operativas:

\begin{enumerate}
  \item ¿Qué casos de uso, fuentes y límites de actuación deben definirse antes
  de configurar el agente?
  \item ¿Qué controles de diseño reducen el riesgo de respuestas no
  fundamentadas, divulgación de información sensible o instrucciones fuera de
  alcance?
  \item ¿Cómo debe evaluarse la calidad de las respuestas de forma repetible,
  separando precisión, cobertura, trazabilidad, consistencia y utilidad?
\end{enumerate}

Estas cuestiones transforman el proyecto en algo más que una demostración de
conversación: el artefacto se diseña junto con un protocolo de evaluación y una
argumentación de seguridad. La metodología elegida se desarrolla en el
Capítulo~\ref{chap:objetivos-metodologia}.

\section{Alcance}
\label{sec:alcance}

El alcance funcional previsto se limita a un escenario organizativo simulado y
a consultas de carácter defensivo y educativo. Las categorías iniciales son las
siguientes:

\begin{itemize}
  \item identificación y comunicación inicial de correos, enlaces o adjuntos
  sospechosos;
  \item gestión de contraseñas, autenticación multifactor y protección de
  cuentas;
  \item pautas básicas para dispositivos, actualizaciones, copias de seguridad y
  uso seguro de servicios corporativos;
  \item orientación inicial ante señales de incidente y derivación al canal de
  notificación; y
  \item explicación de políticas corporativas simuladas y de fuentes públicas
  seleccionadas.
\end{itemize}

La base de conocimiento se compondrá de documentos o sitios oficiales de NIST,
ENISA, INCIBE y Microsoft Learn, complementados, si procede, por una política
corporativa ficticia creada expresamente para las pruebas. Cada fuente deberá
tener propietario, URL o identificador, fecha de revisión, alcance temático y
estado de aprobación. Las fuentes públicas no se tratarán como evidencia
inmutable: deberán revisarse antes de cada versión evaluada, ya que las guías,
las plataformas y el panorama de amenazas evolucionan.

Quedan fuera del alcance los siguientes aspectos:

\begin{itemize}
  \item la monitorización de sistemas, el análisis forense, la ejecución de
  comandos, la respuesta automática ante incidentes o el acceso a herramientas
  de seguridad;
  \item la recepción o tratamiento intencionado de credenciales, secretos,
  registros de producción, datos personales reales o información clasificada;
  \item las instrucciones ofensivas que puedan facilitar intrusión, evasión,
  malware o exfiltración;
  \item la sustitución de un SOC, de un CSIRT, de los responsables de protección
  de datos o de asesoramiento jurídico; y
  \item la afirmación de conformidad legal o certificación de seguridad de una
  organización real.
\end{itemize}

Esta delimitación es deliberada. El agente debe indicar que, ante un riesgo
inmediato o una posible brecha, la persona usuaria ha de seguir el procedimiento
organizativo y contactar con el equipo responsable. Si no existe tal
procedimiento en el entorno de demostración, el texto de derivación se marcará
como un elemento configurable, no como una instrucción universal.

\section{Consideraciones de seguridad, privacidad y gobierno}
\label{sec:consideraciones-intro}

El sistema integra dos superficies que requieren atención: el comportamiento del
agente y la aplicación web que lo presenta. Entre los riesgos característicos de
las aplicaciones con modelos de lenguaje se encuentran la inyección de
instrucciones, la divulgación de información sensible, la contaminación de datos
y la dependencia excesiva de las salidas del modelo \cite{owaspLlm2026}. En
consecuencia, las instrucciones del agente deben priorizar fuentes aprobadas,
prohibir la solicitud de secretos, limitar la respuesta a su dominio y disponer
de rutas de rechazo o escalado.

En el plano web, Angular ofrece mecanismos de protección frente a XSS al tratar
los valores de las plantillas como no confiables y recomienda mantener el
framework actualizado, evitar APIs de riesgo y complementar las defensas con una
política de seguridad de contenido y \emph{Trusted Types} cuando sea aplicable
\cite{angularSecurity}. Estas medidas no sustituyen a la autenticación,
autorización, protección de secretos ni validación del lado servidor; forman
parte de una defensa en profundidad que se concretará durante la implementación.

El Reglamento General de Protección de Datos exige que el tratamiento de datos
personales se realice con medidas técnicas y organizativas adecuadas al riesgo
\cite{gdpr}. Además, el Reglamento europeo de IA establece un marco armonizado
para el desarrollo y uso de sistemas de IA orientado a una IA fiable y centrada
en las personas \cite{aiAct}. Este TFM no emite un dictamen de cumplimiento.
Sí incorpora esos marcos como criterios de diseño: minimización de datos,
transparencia sobre las limitaciones, supervisión humana y documentación de las
decisiones de configuración.

\section{Estructura de la memoria}
\label{sec:estructura-memoria}

El Capítulo~\ref{chap:objetivos-metodologia} define los objetivos, requisitos,
metodología y protocolo de evaluación. El Capítulo~\ref{chap:estado-arte}
revisa los fundamentos de agentes con recuperación de conocimiento, Microsoft
Copilot Studio, seguridad de aplicaciones web y gobierno de IA. Los capítulos
posteriores describirán la arquitectura implementada, la configuración del
agente, la integración con Angular y las decisiones de seguridad. A continuación
se documentará la validación mediante casos de prueba sintéticos y sus resultados
reproducibles. Por último, las conclusiones contrastarán los objetivos con la
evidencia obtenida y propondrán mejoras para un posible despliegue organizativo.

